\documentclass[letterpaper,12pt]{article}
\usepackage[margin=0.55in]{geometry}
\usepackage{tikz}
\usetikzlibrary{arrows.meta}
\usepackage[T1]{fontenc}
\usepackage{lmodern}
\pagestyle{empty}
\setlength{\parindent}{0pt}
\newcommand{\bodyfont}{\fontsize{13.5}{17}\selectfont}
\newcommand{\puttext}[4]{\node[anchor=north west,inner sep=0pt,text width=#3in,align=left] at (#1,-#2) {\bodyfont #4};}
\newcommand{\pagebase}[2]{%
  \node[anchor=north west,inner sep=0pt] at (0,0) {\fontsize{12.5}{15}\selectfont Week 7 / Take-away games / #1};
  \draw[line width=.45pt] (0,-.30) -- (7.4,-.30);
  \node[anchor=west,inner sep=0pt] at (0,-9.82) {\fontsize{9.5}{11}\selectfont Bellingham Math Circle / Week 7 / F07-RV-v1};
  \node[anchor=east,inner sep=0pt] at (7.4,-9.82) {\fontsize{9.5}{11}\selectfont #2};}
\newcommand{\counter}[2]{\fill[black!78] (#1,-#2) circle[radius=.065];}
\newcommand{\passcard}[3]{%
  \draw[rounded corners=2pt,line width=.7pt,fill=white] (#1-.38,-#2-.23) rectangle (#1+.38,-#2+.23);
  \node at (#1,-#2) {\fontsize{10}{11}\selectfont #3};}
\newcommand{\lastcard}[3]{%
  \draw[rounded corners=2pt,line width=.7pt,fill=white] (#1-.39,-#2-.28) rectangle (#1+.39,-#2+.28);
  \node at (#1,-#2+.09) {\fontsize{8}{9}\selectfont LAST};
  \node at (#1,-#2-.10) {\fontsize{13}{14}\selectfont #3};}
\newcommand{\tinyrow}[4]{%
  \foreach \i in {1,...,#3} {
    \pgfmathsetmacro{\xx}{#1+(\i-1)*.20}
    \draw[line width=.45pt] (\xx,-#2) rectangle (\xx+.20,-#2-.20);
    \ifnum\i=#4\relax\else\fill[black!78] (\xx+.10,-#2-.10) circle[radius=.055];\fi
  }}
\newcommand{\board}[3]{%
  \foreach \i in {1,...,#3} {
    \pgfmathsetmacro{\xx}{#1+(\i-1)*.80}
    \draw[line width=.7pt] (\xx,-#2) rectangle (\xx+.80,-#2-.80);
  }}
\begin{document}
\null
\begin{tikzpicture}[remember picture,overlay,x=1in,y=1in,>=Stealth]
\begin{scope}[shift={([xshift=.55in,yshift=-.55in]current page.north west)}]
\pagebase{Grades 2--5}{1}
\puttext{0}{.55}{7.4}{Take turns taking 1 or 2 counters. Either player may turn the shared face-up PASS card face down instead of taking counters. Taking the last counter ends the game and wins.}

\foreach \x/\label in {1.10/Before,3.70/Turn the card over,6.30/After} {
  \node at (\x,-1.65) {\fontsize{11}{12}\selectfont \label};
  \node at (\x,-1.91) {\fontsize{10}{11}\selectfont 14 counters};
  \foreach \j in {0,...,6} {
    \pgfmathsetmacro{\cx}{\x-.66+\j*.22}
    \counter{\cx}{2.16}\counter{\cx}{2.40}
  }
}
\passcard{1.10}{2.84}{PASS}
\begin{scope}[shift={(3.70,-2.84)},rotate=32]
  \draw[rounded corners=2pt,line width=.7pt,fill=black!8] (-.38,-.23) rectangle (.38,.23);
  \node {\fontsize{10}{11}\selectfont PASS};
\end{scope}
\draw[->,line width=.65pt] (3.24,-2.75) to[out=95,in=85] (4.15,-2.78);
\passcard{6.30}{2.84}{used}
\draw[->,line width=.7pt] (2.18,-2.31) -- (2.71,-2.31);
\draw[->,line width=.7pt] (4.71,-2.31) -- (5.20,-2.31);

\puttext{0}{3.40}{7.4}{\textbf{Problem 1:} Play with 1 to 12 counters. When does one PASS card change which player can make sure they win? Circle the starts where the first player can make sure they win in each list.}
\draw[rounded corners=3pt,line width=.45pt] (0,-4.65) rectangle (3.55,-7.35);
\draw[rounded corners=3pt,line width=.45pt] (3.85,-4.65) rectangle (7.40,-7.35);
\node at (1.775,-4.98) {\bodyfont PASS face up};
\node at (5.625,-4.98) {\bodyfont PASS already used};
\foreach \offset in {0,3.85} {
  \foreach \n in {1,...,12} {
    \pgfmathtruncatemacro{\rownum}{int((\n-1)/4)}
    \pgfmathtruncatemacro{\colnum}{mod(\n-1,4)}
    \pgfmathsetmacro{\xx}{\offset+.47+.87*\colnum}
    \pgfmathsetmacro{\yy}{5.58+.64*\rownum}
    \node at (\xx,-\yy) {\fontsize{18}{21}\selectfont \n};
  }
}
\end{scope}
\end{tikzpicture}

\newpage\null
\begin{tikzpicture}[remember picture,overlay,x=1in,y=1in,>=Stealth]
\begin{scope}[shift={([xshift=.55in,yshift=-.55in]current page.north west)}]
\pagebase{Grades 4--5}{2}
\puttext{0}{.55}{7.4}{Take turns taking 1, 2 or 3 counters. You may not take the number on the face-up LAST card. After each turn, turn the old LAST card face down and the matching card face up. Taking the last counter wins; if you cannot make a legal move, you lose. No face-up LAST card means there is no previous move.}

\foreach \x/\label in {1.10/Before,3.70/Take 2,6.30/After} {
  \node at (\x,-1.96) {\fontsize{11}{12}\selectfont \label};
}
\node at (1.10,-2.21) {\fontsize{10}{11}\selectfont 13 counters};
\node at (3.70,-2.21) {\fontsize{10}{11}\selectfont remove 2};
\node at (6.30,-2.21) {\fontsize{10}{11}\selectfont 11 counters};
\foreach \j in {0,...,6} {
  \pgfmathsetmacro{\cx}{.44+\j*.22}
  \counter{\cx}{2.45}
}
\foreach \j in {0,...,5} {
  \pgfmathsetmacro{\cx}{.44+\j*.22}
  \counter{\cx}{2.69}
}
\foreach \j in {0,...,6} {
  \pgfmathsetmacro{\cx}{3.04+\j*.22}
  \counter{\cx}{2.45}
}
\foreach \j in {0,...,3} {
  \pgfmathsetmacro{\cx}{3.04+\j*.22}
  \counter{\cx}{2.69}
}
\foreach \j in {4,5} {
  \pgfmathsetmacro{\cx}{3.04+\j*.22}
  \fill[black!15] (\cx,-2.69) circle[radius=.065];
  \draw[line width=.7pt] (\cx-.05,-2.64) -- (\cx+.05,-2.74);
  \draw[line width=.7pt] (\cx-.05,-2.74) -- (\cx+.05,-2.64);
}
\foreach \j in {0,...,6} {
  \pgfmathsetmacro{\cx}{5.64+\j*.22}
  \counter{\cx}{2.45}
}
\foreach \j in {0,...,3} {
  \pgfmathsetmacro{\cx}{5.64+\j*.22}
  \counter{\cx}{2.69}
}
\lastcard{1.10}{3.09}{1}
\lastcard{3.70}{3.09}{1}
\lastcard{6.30}{3.09}{2}
\draw[->,line width=.7pt] (2.18,-2.58) -- (2.71,-2.58);
\draw[->,line width=.7pt] (4.71,-2.58) -- (5.20,-2.58);

\puttext{0}{3.72}{7.4}{\textbf{Problem 2:} Can changing the LAST card change whether the first player can make sure they win? Play with piles of 1 to 8 counters and circle the winning starts in each column. Find a rule that predicts the winners for larger piles.}
\draw[line width=.45pt] (0,-4.92) rectangle (7.40,-8.68);
\foreach \x in {1.85,3.70,5.55} {\draw[line width=.45pt] (\x,-4.92) -- (\x,-8.68);}
\node[align=center] at (.925,-5.28) {\fontsize{11}{12}\selectfont No previous\\move};
\lastcard{2.775}{5.28}{1}
\lastcard{4.625}{5.28}{2}
\lastcard{6.475}{5.28}{3}
\draw[line width=.45pt] (0,-5.67) -- (7.40,-5.67);
\foreach \n in {1,...,8} {
  \pgfmathsetmacro{\yy}{5.86+(\n-1)*.375}
  \foreach \x in {.925,2.775,4.625,6.475} {
    \node at (\x,-\yy) {\fontsize{16}{18}\selectfont \n};
  }
}
\draw[black!35,line width=.4pt] (0,-9.02) -- (7.40,-9.02);
\draw[black!35,line width=.4pt] (0,-9.44) -- (7.40,-9.44);
\end{scope}
\end{tikzpicture}

\newpage\null
\begin{tikzpicture}[remember picture,overlay,x=1in,y=1in,>=Stealth]
\begin{scope}[shift={([xshift=.55in,yshift=-.55in]current page.north west)}]
\pagebase{Grades K--5}{3}
\puttext{0}{.55}{7.4}{Take turns removing exactly 2 counters from neighbouring squares in one row. Gaps stay empty; counters do not slide. If you cannot take a pair, you lose.}
\foreach \x/\label in {.10/Before,2.98/Take this pair,5.90/After} {
  \node[anchor=west] at (\x,-1.58) {\fontsize{10}{11}\selectfont \label};
}
\foreach \i in {1,...,7} {
  \pgfmathsetmacro{\xx}{.10+(\i-1)*.20}
  \draw[line width=.45pt] (\xx,-1.85) rectangle (\xx+.20,-2.05);
  \fill[black!78] (\xx+.10,-1.95) circle[radius=.055];
  \pgfmathsetmacro{\xx}{2.98+(\i-1)*.20}
  \draw[line width=.45pt] (\xx,-1.85) rectangle (\xx+.20,-2.05);
  \ifnum\i=2
    \fill[black!15] (\xx+.10,-1.95) circle[radius=.055];
    \draw[line width=.7pt] (\xx+.05,-1.90) -- (\xx+.15,-2.00);
    \draw[line width=.7pt] (\xx+.05,-2.00) -- (\xx+.15,-1.90);
  \else\ifnum\i=3
    \fill[black!15] (\xx+.10,-1.95) circle[radius=.055];
    \draw[line width=.7pt] (\xx+.05,-1.90) -- (\xx+.15,-2.00);
    \draw[line width=.7pt] (\xx+.05,-2.00) -- (\xx+.15,-1.90);
  \else\fill[black!78] (\xx+.10,-1.95) circle[radius=.055];\fi\fi
  \pgfmathsetmacro{\xx}{5.90+(\i-1)*.20}
  \draw[line width=.45pt] (\xx,-1.85) rectangle (\xx+.20,-2.05);
  \ifnum\i=2\else\ifnum\i=3\else\fill[black!78] (\xx+.10,-1.95) circle[radius=.055];\fi\fi
}
\draw[->,line width=.7pt] (1.67,-1.95) -- (2.76,-1.95);
\draw[->,line width=.7pt] (4.55,-1.95) -- (5.67,-1.95);

\puttext{0}{2.42}{7.4}{\textbf{Problem 3:} Choose a start and put a counter in each square. Play all six starts and circle those where the first player can make sure they win.}

\node at (1.60,-3.72) {\bodyfont A};
\board{0}{4.10}{4}
\node at (5.80,-3.45) {\bodyfont B};
\board{5.00}{3.70}{2}
\board{5.00}{4.60}{2}
\draw[line width=.65pt] (4.87,-3.70) -- (4.77,-3.70) -- (4.77,-5.40) -- (4.87,-5.40);

\node at (2.00,-5.60) {\bodyfont C};
\board{0}{5.85}{5}
\node at (6.00,-5.60) {\bodyfont D};
\board{4.80}{5.85}{3}

\node at (3.00,-6.90) {\bodyfont E};
\board{.60}{7.15}{6}
\node at (3.70,-8.20) {\bodyfont F};
\board{.10}{8.45}{9}
\end{scope}
\end{tikzpicture}
\end{document}
